% answers_short.tex —— 2004 年·普通高等学校招生考试·江苏卷 · 简短答案速查
% 由 tools/make_short_answers.py 自 answers.tex 的 \gksolhead/\gkanswer 提取，
% 并与 2004/2004_数学真题_江苏_答案要点.md 逐题比对后生成，勿手改（改源 MD 后重跑）。

\begin{gknotice}
\textbf{使用说明}：本节只给最终答案便于快速核对，完整推导见解析版。答案为编写者独立推导，\textbf{非官方标准答案}。
\end{gknotice}

\gkpart{选择题}
\begin{center}\begin{tabular}{ccccccccccccc}
\toprule
题号 & 1 & 2 & 3 & 4 & 5 & 6 & 7 & 8 & 9 & 10 & 11 & 12 \\
\midrule
答案 & A & B & D & C & A & B & C & A & D & C & B & A \\
\bottomrule
\end{tabular}\end{center}

\gkpart{填空题}
\begin{itemize}[leftmargin=2.6em]
  \item[\textbf{13.}] $\{x \mid x < -2 \text{ \text{或} } x > 3\}$（或 $(-\infty, -2) \cup (3, +\infty)$）
  \item[\textbf{14.}] $(x - 1)^2 + (y - 2)^2 = 25$
  \item[\textbf{15.}] $2$
  \item[\textbf{16.}] $\left(\dfrac{4}{5}, -\dfrac{3}{5}\right)$
\end{itemize}

\gkpart{解答题}
\begin{itemize}[leftmargin=2.6em]
  \item[\textbf{17.}] $\sin\left(\alpha - \dfrac{\pi}{3}\right) = \dfrac{4 - 3\sqrt{3}}{10}$
  \item[\textbf{18.}] （1）直线 $AP$ 与平面 $BCC_1B_1$ 所成的角的大小为 $\arctan\dfrac{4\sqrt{17}}{17}$（或 $\arctan\dfrac{4}{\sqrt{17}}$）；（2）证明见解析；（3）点 $P$ 到平面 $ABD_1$ 的距离为 $\dfrac{3\sqrt{2}}{2}$
  \item[\textbf{19.}] 投资人对甲项目投资 4 万元，对乙项目投资 6 万元，才能使可能的盈利最大，最大盈利为 7 万元
  \item[\textbf{20.}] （1）$k = 4$；（2）满足条件的无穷等差数列 $\{a_n\}$ 共有三个：$a_n = 0$；$a_n = 1$；$a_n = 2n - 1$
  \item[\textbf{21.}] （1）椭圆的方程为 $\dfrac{x^2}{4m^2} + \dfrac{y^2}{3m^2} = 1$；（2）直线 $l$ 的斜率为 $\pm 2\sqrt{6}$
  \item[\textbf{22.}] （1）证明见解析；（2）证明见解析；（3）证明见解析
\end{itemize}
